\documentclass[11pt]{amsart}

\usepackage[T1]{fontenc}
\usepackage{lmodern}
\usepackage{microtype}
\usepackage{mathtools,amssymb,amsthm}
\usepackage[hidelinks]{hyperref}

\hypersetup{
  pdftitle={Counterexamples to the Zhang--Wan Deep-Hole Completeness Conjecture}
}

\newtheorem{theorem}{Theorem}
\newtheorem{lemma}[theorem]{Lemma}
\newtheorem{proposition}[theorem]{Proposition}

\newcommand{\F}{\mathbb{F}}
\newcommand{\CL}{C_{\mathcal L}}
\newcommand{\PG}{\operatorname{PG}}

\title[Counterexamples to the Zhang--Wan conjecture]
{Counterexamples to the Zhang--Wan Deep-Hole Completeness Conjecture}
\date{}
\subjclass[2020]{Primary 94B27; Secondary 14G50}
\keywords{elliptic curve code, deep hole, covering radius, syndrome}

\begin{document}

\begin{abstract}
Let $E/\F_q$ be an elliptic curve of odd characteristic, let
$N=\#E(\F_q)\ge q+4$, and put $D=E(\F_q)\setminus\{O\}$.
Zhang and Wan conjectured that, for $2\le k\le N-4$, the layer
$\CL(D,(k+1)O)\setminus\CL(D,kO)$ is the complete set of deep holes of
$\CL(D,kO)$. We refute the conjecture as stated. At the upper endpoint
$k=N-4$, the code has covering radius two and
$(q-1)(q^2+q+2-N)$ deep-hole cosets, while the proposed layer contains
only $q-1$ cosets. We also give an example over $\F_3$ for which the
conjecture fails at both admissible values $k=2$ and $k=3$.
\end{abstract}

\maketitle

For a linear code $C\subseteq\F_q^n$, let $\rho(C)$ denote its covering
radius. A word $v\in\F_q^n$ is a \emph{deep hole} if
$d(v,C)=\rho(C)$. Distance to $C$ is constant on cosets of $C$.

Zhang and Wan \cite[Conjecture~1.4]{ZhangWan} conjectured the following.
Let $E/\F_q$ be an elliptic curve, where $q$ is odd, with
$N=\#E(\F_q)\ge q+4$. For $O\in E(\F_q)$ and
$D=E(\F_q)\setminus\{O\}$, they asserted that
\begin{equation}\label{eq:ZW}
  \CL(D,(k+1)O)\setminus\CL(D,kO)
\end{equation}
is the set of all deep holes of $\CL(D,kO)$ whenever
$2\le k\le N-4$. We first show that the upper endpoint always fails.

\section{The endpoint obstruction}

\begin{lemma}\label{lem:codim-three}
Let $C$ be an $[n,n-3]$ code over $\F_q$ with $d(C)\ge3$. If
\[
  q+1<n<q^2+q+1,
\]
then $\rho(C)=2$, and the number of deep-hole cosets of $C$ is
\[
  q^3-1-n(q-1)=(q-1)(q^2+q+1-n).
\]
\end{lemma}

\begin{proof}
Let $H$ be a $3\times n$ parity-check matrix. Since $d(C)\ge3$, the
columns of $H$ are nonzero and pairwise nonproportional, and hence define
$n$ distinct points of $\PG(2,q)$.

Let $s\in\F_q^3\setminus\{0\}$. If $[s]$ is a column point, then $s$ is
generated by one column. Otherwise, the $q+1$ projective lines through
$[s]$ partition the column points. Since $n>q+1$, one of these lines
contains two column points, whose representatives span $s$. Thus every
syndrome is generated by at most two columns, so $\rho(C)\le2$.

Since $\PG(2,q)$ has $q^2+q+1$ points and $n<q^2+q+1$, some nonzero
syndrome is not proportional to a column. Its coset weight is two, and
therefore $\rho(C)=2$. Of the $q^3$ syndromes, one has weight zero and
exactly $n(q-1)$ have weight one. Every remaining syndrome has weight
two, giving the stated count.
\end{proof}

\begin{theorem}\label{thm:endpoint}
Let $q$ be an odd prime power, let $E/\F_q$ be an elliptic curve with
$N=\#E(\F_q)\ge q+4$, and fix $O\in E(\F_q)$. Put
\[
  D=E(\F_q)\setminus\{O\},\qquad n=N-1,\qquad k=N-4=n-3,
\]
and let $C=\CL(D,kO)$. Then $\rho(C)=2$, and $C$ has
\[
  (q-1)(q^2+q+1-n)
\]
deep-hole cosets. The layer $\CL(D,(k+1)O)\setminus C$ is the union of
exactly $q-1$ deep-hole cosets and is a proper subset of the deep holes
of $C$. Consequently, Conjecture~1.4 of Zhang and Wan is false as stated;
its hypotheses are satisfiable, as Proposition~\ref{prop:F3} shows.
\end{theorem}

\begin{proof}
Since $E$ has genus one, Riemann--Roch gives $\ell(kO)=k$ for $k\ge1$
(here $k=N-4\ge q\ge3$). Evaluation on
$D$ is injective because $\deg(kO-D)=k-n<0$. Hence $C$ has dimension
$k=n-3$, and the standard designed-distance bound gives $d(C)\ge n-k=3$;
see \cite[Chapter~2]{Stichtenoth}.

The hypothesis $N\ge q+4$ and Hasse's bound give
\[
  q+3\le n\le q+2\sqrt q.
\]
Thus $q+1<n<q^2+q+1$, and Lemma~\ref{lem:codim-three} gives the covering
radius and the number of deep-hole cosets.

Set $C'=\CL(D,(k+1)O)$. Again by Riemann--Roch and injectivity of
evaluation,
\[
  \dim C'=k+1,\qquad d(C')\ge n-(k+1)=2.
\]
Thus $C'/C$ is one-dimensional. Each of its $q-1$ nonzero cosets has
coset weight at least two, since $C'$ contains no word of weight one.
As $\rho(C)=2$, these are deep-hole cosets.

Finally,
\[
  n\le q+2\sqrt q<q^2+q,
\]
so $q^2+q+1-n>1$. Hence the total number of deep-hole cosets is strictly
larger than $q-1$, and the layer is not complete.
\end{proof}

Two of the assertions of Theorem~\ref{thm:endpoint} are already those of
the source. Since $\#E(\F_q)-n=1$ for $D=E(\F_q)\setminus\{O\}$, the
count $(\#E(\F_q)-n)(q-1)q^{k}$ of \cite[Theorem~1.2(iii)--(iv)]{ZhangWan}
specialises to $(q-1)q^{k}$ words, that is, to the $q-1$ cosets of the
layer, and the hypothesis $k<n-2$ of their part~(iv) holds at $k=n-3$;
likewise $\rho(C)=n-k-1=2$ is their part~(ii). What Theorem~\ref{thm:endpoint}
adds is, first, that both hold with no side condition: their theorem
requires, besides $n\ge q+3$, one of $n\ge q+k$, $q$ prime, or
$k\le\sqrt q$, and at $k=n-3$
the first forces $q\le3$ while the third fails because $k\ge q$, so for
non-prime $q$ their argument does not reach the endpoint, whereas
Lemma~\ref{lem:codim-three} gives $\rho(C)=2$ unconditionally. Second,
and this is the point, it supplies the total number
$(q-1)(q^2+q+1-n)$ of deep-hole cosets, against which the layer can be
compared.

The cap $k\le N-4$ is deliberate: Zhang and Wan note in their
Remark~1.3 that at $k=n-2=N-3$ the covering radius drops to $n-k-1=1$
(again under one of those conditions), so every word outside $C$ is a
deep hole and completeness fails there for a trivial reason.
Theorem~\ref{thm:endpoint} shows that it already fails at $k=N-4$, the
largest value the conjecture does admit, and Proposition~\ref{prop:F3}
shows it can fail at the smallest admissible value $k=2$ as well; the
failure is therefore not an artefact of the endpoint. We note finally
that \cite[Corollary~4.11]{ZhangWan} does prove a conditional
completeness statement in codimension three, but for the residue codes
$C_\Omega(D,kO)$ rather than for the functional codes $\CL(D,kO)$ of
Conjecture~1.4; it does not conflict with Theorem~\ref{thm:endpoint}.

\section{Failure at both parameters over \texorpdfstring{$\F_3$}{F3}}

\begin{proposition}\label{prop:F3}
Over $\F_3$, let
\[
  E:\ y^2=x^3+2x+1,
\]
let $O$ be the point at infinity, and put
$D=E(\F_3)\setminus\{O\}$. Then $\#E(\F_3)=7$, and the Zhang--Wan
conjecture fails for both admissible parameters $k=2$ and $k=3$.
\end{proposition}

\begin{proof}
The projective closure is
\[
  F(X,Y,Z)=Y^2Z-X^3-2XZ^2-Z^3=0.
\]
Over $\F_3$, one has $F_X=Z^2$, so a singular point would have $Z=0$.
The unique point at infinity is $O=[0:1:0]$, and $F_Z(O)=1$; hence $E$ is
nonsingular. Moreover, $x^3+2x+1=1$ for every $x\in\F_3$, so
$\#E(\F_3)=7$. Order the points of $D$ as
\[
  (0,1),\ (0,2),\ (1,1),\ (1,2),\ (2,1),\ (2,2);
\]
all coordinates below refer to this order.
The functions $x$ and $y$ have their only poles at
$O$, of orders two and three, respectively. Hence Riemann--Roch gives
\[
  L(2O)=\langle1,x\rangle,\qquad L(3O)=\langle1,x,y\rangle.
\]

First take $k=2$ and write $C_2=\CL(D,2O)$. Then
\begin{align*}
 C_2
 &=\{(a,a,a+b,a+b,a+2b,a+2b):a,b\in\F_3\}\\
 &=\left\{(u_0,u_0,u_1,u_1,u_2,u_2):
 u_0+u_1+u_2=0\right\}.
\end{align*}
View a received word as three consecutive coordinate pairs. Every word
of $\F_3^6$ agrees with a word of $C_2$ in at least three coordinates.
Indeed, if one pair is constant, match both entries there, match one
entry in a second pair, and choose the value on the third pair to make
$u_0+u_1+u_2=0$. If all three pairs are nonconstant, let
$A_i\subset\F_3$ be the two symbols occurring in the $i$th pair. Choose
$u_2\in A_2$. Since any two two-element subsets of $\F_3$ satisfy
$A_0+A_1=\F_3$, there are $u_0\in A_0$ and $u_1\in A_1$ with
$u_0+u_1=-u_2$, matching one entry in every pair. Thus $\rho(C_2)\le3$.

Now let
\[
  w=(0,1,0,1,0,2).
\]
Every pair of $w$ is nonconstant, so every word of $C_2$ differs from
$w$ in at least one coordinate of each pair; hence $d(w,C_2)\ge3$.
Since $\operatorname{wt}(w)=3$ and $0\in C_2$, this lower bound is
attained, so $d(w,C_2)=3$. Hence $\rho(C_2)=3$ and $w$ is a deep hole.
For $f=a+bx+cy\in L(3O)$, the difference
$f(x,1)-f(x,2)=-c$ is independent of $x$. The three corresponding pair
differences of $w$ are $2,2,1$, so $w\notin\CL(D,3O)$. Therefore $w$ is
a deep hole of $C_2$ outside the proposed layer
$\CL(D,3O)\setminus C_2$. Since $\dim\CL(D,3O)=3$, that layer is the
union of exactly two cosets of $C_2$, and both are deep-hole cosets:
every nonzero word of $\CL(D,3O)$ has weight at least $6-3=3$, so each of
the two has coset weight $3=\rho(C_2)$. Thus at $k=2$ the layer is made
of deep holes and is still incomplete, $w$ exhibiting a third deep-hole
coset. A direct enumeration of $\F_3^6$, on which nothing here depends,
shows that $24$ of the $81$ cosets of $C_2$ have coset weight three, so
that $22$ of them lie outside the layer.

For $k=3=N-4$, Theorem~\ref{thm:endpoint} applies. The code
$\CL(D,3O)$ has fourteen deep-hole cosets, whereas
$\CL(D,4O)\setminus\CL(D,3O)$ consists of only two cosets. Thus the
conjecture also fails at $k=3$.
\end{proof}

\medskip
\noindent\textbf{Exact verification.}
The one finite computation quoted above is the coset census at the end of the
proof of Proposition~\ref{prop:F3}, and it is recomputable from the data
printed in this paper. From $L(2O)=\langle1,x\rangle$ and the point order fixed
in that proof one writes down the nine words of $C_2$, forms the $81$ cosets of
$C_2$ in $\F_3^6$, and takes the minimum weight in each: exactly $24$ have
coset weight three, and since the layer $\CL(D,3O)\setminus C_2$ accounts for
two of these, $22$ of them lie outside the layer. The same printed data carry
the two checks the argument actually uses, both of them by hand:
$d(w,C_2)=3$ for the word $w$ displayed above, and $w\notin\CL(D,3O)$ because
the three pair differences of $w$ are not constant. The count of fourteen
deep-hole cosets of $\CL(D,3O)$ is Lemma~\ref{lem:codim-three} evaluated at
$q=3$ and $n=6$, not a computation. All arithmetic is exact, in $\F_3$ and on
integers; no floating-point computation is involved.

The census was also re-run independently, by a program written from the
definitions of this note rather than from the first implementation: its $81$
coset-weight determinations, and the two counts $24$ and $22$ derived from
them, agreed throughout with the values printed above. Both programs and the
transcript of that run are retained in an auxiliary archive and are available
on request; the specification above is what a reader should reimplement. A
program accompanying this note carries out that reimplementation: from the
field, the curve, the point order and the word printed above it re-derives
$C_2$, the $81$ coset weights and the counts $24$ and $22$, and also the
covering radius and the deep-hole count of Theorem~\ref{thm:endpoint} for this
curve and for a bounded family of curves over small prime fields. That archive
is separate from it and is not distributed with this note. No
part of Theorem~\ref{thm:endpoint} or of Proposition~\ref{prop:F3} depends on
the census, which is reported only to quantify how far the layer falls short
at $k=2$.

\section*{Status and scope}

Theorem~\ref{thm:endpoint} refutes Conjecture~1.4 of \cite{ZhangWan} as
stated: for every odd prime power $q$, every elliptic curve $E/\F_q$ with
$\#E(\F_q)\ge q+4$ and every $O\in E(\F_q)$, the conjecture fails at
$k=N-4$; Proposition~\ref{prop:F3} exhibits an instance in which it also
fails at $k=2$. What is not settled here is the interior of the range.
For $q>3$ we do not decide, for any single $k<N-4$, whether
$\CL(D,(k+1)O)\setminus\CL(D,kO)$ is the complete set of deep holes of
$\CL(D,kO)$, and we do not propose a corrected range. In both cases
treated here the proposed layer does consist of deep holes, so what fails
is completeness alone.

The conjecture is quoted from the preprint version of \cite{ZhangWan},
whose statement numbering we follow; we were not able to consult the
printed version, in which the numbering may differ. We are aware of no
resolution of Conjecture~1.4 in the literature and of no earlier
appearance of the word $w$ of Proposition~\ref{prop:F3}; that is a
bounded negative obtained from the citing literature and not a proof of
absence. No novelty is claimed for Lemma~\ref{lem:codim-three}, which is
a standard syndrome count in $\PG(2,q)$; the same geometric viewpoint is
used in \cite[Proposition~4.9]{ZhangWan}.

\begin{thebibliography}{9}

\bibitem{Stichtenoth}
H.~Stichtenoth,
\emph{Algebraic Function Fields and Codes}, 2nd ed.,
Graduate Texts in Mathematics, vol.~254,
Springer, Berlin, 2009.

\bibitem{ZhangWan}
J.~Zhang and D.~Wan,
\emph{On deep holes of elliptic curve codes},
IEEE Trans. Inform. Theory \textbf{69} (2023), no.~7, 4498--4506,
\href{https://doi.org/10.1109/TIT.2023.3257320}
{doi:10.1109/TIT.2023.3257320};
preprint \href{https://arxiv.org/abs/2207.12584}{arXiv:2207.12584}.
Statement numbers cited above are those of the preprint, in which the
completeness conjecture is Conjecture~1.4, the covering-radius and
deep-hole theorem is Theorem~1.2, the boundary case $k=n-2$ is discussed
in Remark~1.3, and the conditional completeness result for residue codes
is Corollary~4.11.

\end{thebibliography}

\end{document}